\documentclass{article}

\usepackage{amsfonts} % Common number sets using \mathbb 
\usepackage{amsthm} % theorems
\usepackage{amsmath} % \implies

\usepackage[utf8]{inputenc}
\usepackage[T1]{fontenc}

\usepackage{bookmark} % used for having pdf bookmarks

\title{Nate's Basic Concepts Notes}
\author{Nate Harris}

\theoremstyle{definition}
\newtheorem{thm}{Theorem}[section]

\theoremstyle{definition}
\newtheorem{dfn}[thm]{Definition}
\newtheorem{stmt}[thm]{Statement}
\newtheorem{exmp}[thm]{Example}

\AfterEndEnvironment{thm}{\noindent\ignorespaces}
\AfterEndEnvironment{dfn}{\noindent\ignorespaces}
\AfterEndEnvironment{exmp}{\noindent\ignorespaces}

\begin{document}
\maketitle{}
\newpage
\tableofcontents{}
\newpage

\section*{Chapter 1}
\addcontentsline{toc}{section}{Chapter 1}

We're starting at the 3rd section of the chapter... first 2 just kinda tell you
why we're doing this and to have fun! And also to be organized and write proofs in
full sentences that are easily readable and hard to misinterpret.
\setcounter{section}{3}
\section*{Definitions}
\addcontentsline{toc}{subsection}{Definitions}

Making definitions requires that the language is interpretable in only one way.
The basic understandings later definitions will be built upon are... 
We know how to add, subtract, multiply and the algebraic properties related.
We also understand \textgreater, \textless, \underline{\textless}, \underline{\textgreater}
\\[\baselineskip]
\noindent We also define integers here as:
\begin{center}
    $\mathbb{Z} = \{..., -3, -2, -1, 0, 1, 2, 3, ...\}$.
\end{center}
\begin{dfn} \textbf{Even Numbers}
    An integer \textit{a} is even if it is divisible by 2.
\end{dfn}

This definition is a little confusing as we haven't defined what it means to 
be divisible by something.

\begin{dfn} \textbf{Divisible}
An integer \textit{a} is divisible by \textit{b} if there is an integer $c$
    such that $bc = a$. If this is true we say $b$ divides $a$, $b$ is a factor of
    $a$ or $b$ is a divisor of $a$. The notation of this is $b|a$.
\end{dfn}

\begin{exmp} \textbf{Is 12 divisible by 4?}
    Yes! 4|12! We can find a c of 3 which satisfies the definition. But... is 13 divisible
    by 4? No... as no integer c can be found that satisfies the definition.
\end{exmp}

Now we have 2 valid definitions as 3.2 ensures all terms used in 3.1 are part of
the expected understanding.
\begin{dfn}\textbf{Odd}
An integer $a$ is odd if there exists some integer $b$ such that $2b + 1 = a $ 
\end{dfn}
\begin{dfn}\textbf{Prime} 
    An integer $p$ is prime if $p < 1$ and the only positive dividers of $p$ are
    $p$ and $1$.
\end{dfn}
    
A little note attached to this is a discussion on why 1 is not prime. And it kind
of boils down to that it is just how the math community established and understands prime
numbers. If you wanted a prime number definition to include one, you would have to
put a word like \textit{relaxed} or something in front of \textit{prime}.

Also there is mention of a prime factorization, which is breaking down the number into
factors that are all prime. i.e. $12 = 1\times2\times2\times3$
\begin{dfn}\textbf{Composite}
    An integer $a$ is considered composite if there is a positive integer $b$ s.t.
    $1 < b < a$ and $b|a$.
\end{dfn}

It is impossible for a prime number to be composite as the only factors of a prime
number $p$ are 1 and $p$ which does not satisfy the inequality in the composite
definition


\section*{Examples}
\addcontentsline{toc}{subsection}{Examples}
The examples here were done in my blue notebook. Just a bit easier
An important thing from them was the textbooks definition of the set of natural
numbers.
\\[\baselineskip ]
The set of Natural numbers can be defined as:
\begin{center}
$\mathbb{N} = \{0,1,2,3,...\}.$
\end{center}
Another important set defined here as well is the set of rational numbers ($\mathbb{Q}$).
\\[\baselineskip]
A rational number is a number that can be represented by two integers $a/b$ where
$b \not= 0$ which is denoted by $\mathbb{Q}$.

One is also told to define a \textit{perfect square}:
\begin{dfn} \textbf{Perfect Square}
   An integer $p$ can be a perfect square if there is an integer $x$ s.t. $x \cdot x = p$
   This operation of multiplying an integer $x$ by itself can also be expressed as $x^2$
   We say $x$ 
\end{dfn}

Also one is asked to define a square root.
\begin{dfn}\textbf{Square Root of a Number}
   A natural number $a$ is the square root of natural number $b$ if $a^2 = b$. This 
   can also be denoted as $a = \sqrt{b}$
\end{dfn}

Asked to define a perimeter.
\begin{dfn}\textbf{Perimeter of Polygon}
    The perimeter $P$ of a polygon $p$ is the summation of the lengths of the sides of 
    this polygon $p$.
\end{dfn}
\begin{dfn}\textbf{Between 2 Points}
    Suppose points on the plane $A$,$B$ and $C$ are points on a plane. We say that $C$ is between $A$ and
    $B$ provided that a straight line segment constructed between points $A$ and $B$ also contains $C$. 
\end{dfn}

\begin{dfn}\textbf{Collinear}
    Points on the plane $A$,$B$ and $C$ are collinear if there is some arrangement of these points such that 
    one of the points is between the other two.
\end{dfn}
\begin{dfn}\textbf{Midpoint}
    Points on the plane $A$,$B$ and $C$ are collinear. $C$ is between $A$ and $B$. $C$ is the midpoint
    if the distances between points $A$ and $C$ and $B$ and $C$ are the same.
\end{dfn}
\setcounter{section}{4}
\setcounter{thm}{0}
\pagebreak
\section*{Theorem}
\addcontentsline{toc}{subsection}{Theorem}
\subsection*{Definition of a Theorem}
A theorem is a declarative statement in math that is 
supported by a proof
\subsubsection*{Definition of a Proof}
A proof is basically a mathematical essay with a conclusion
that shows something is true.
\\[\baselineskip]
There are three kinds of declarative statements in math,
\begin{itemize}
    \item{\textbf{Theorem}: something that is proven to be true}
    \item{\textbf{Conjecture}: something that has yet to or cannot be proven}
    \item{\textbf{Mistakes}: statements that just aren't true. 1+1=4, or the idea
        of a geocentric model of the solar system}
\end{itemize}

A side note here, monitor your language when writing proofs! Don't write \textit{nonsense!}
\subsubsection*{Truth in Math}
An example of a statement that is not strictly true enough for Mathematical standards is
the statement that an object dropped near the surface of the earth accelerates at $9.8m/s^2$.
\textit{Near} is not a clear enough term and 9.8 isn't precisely accurate. Even
if \textit{near} is defined as for example: ``within 100 meters'' it will vary.


Regardless, this statement is true enough for most uses. Though, in mathematics one
deals in universal truth, something that must always be true in all cases and for this to
be the case mathematics doesn't deal in real world truth.

Here is an example of a theorem we all know...
\begin{thm}\textbf{Pythagorean}
    If $a$ and $b$ are lengths of the legs of a right triangle and $c$ is the 
    length of the hypotenuse: 
    \begin{center}
        $a^2 + b^2 = c^2$
    \end{center}
\end{thm}
This is true in all cases. This can be proven which will be done later.
\subsubsection*{If-Then Contraction}
Mathematical language follows different rules than common English.
The most important example of this is the If-Then statement construction.
Many theorems can be expressed in the format ``If $A$, then $B$''. For example:

\begin{center}
    The sum of two even integers is even 
\end{center}
becomes:

\begin{center}
    \textbf{If} $x$ and $y$ are even integers \textbf{then} $x + y$ is even. 
\end{center}
In mathematics the use of if-then only promises one outcome. For example the earlier
example only says that if $x$ and $y$ are even integers then $x + y$ is even. It
doesn't say that if $x$ and $y$ are not even integers then $x+y$ is not even.
The only condition that is promised is that if \textbf{A} is true \textbf{B} is also
true.
\bigbreak


\noindent This table shows what outcomes an if-then statement asserts


\begin{center}
\begin{tabular}{l|l||l}
    A & B & \\
    \hline
    True & True & Possible\\
    True & False & Not Possible\\
    False & True & Possible\\
    False & False & Possible\\
\end{tabular}
\end{center}

Alternate wordings for If \textit{A} then \textit{B}:
\begin{itemize}
    \item``\textit{A} implies \textit{B}'' or ``\textit{B}
        implies \textit{A}''
    \item ``Whenever \textit{A}, we have \textit{B}''
    \item \textit{A} is sufficient for \textit{B}
    \item $A \implies B$
    \item $A \impliedby B$
\end{itemize}
\subsubsection*{If and Only If}
Some theorems can be expressed in the form If $A$ then $B$ and If $B$ then $A$.
For example:
\begin{center}
    If an integer $x$ is even, then $x+1$ is odd and if $x+1$ is odd, then $x$ is even
\end{center}
This can be expressed shorter in If and Only If format.
\begin{center}
    An integer $x$ is even, if and only if $x+1$ is odd.
\end{center}

Here is a truth table that shows what If and Only If statements assert:

\begin{center}
\begin{tabular}{l|l||l}
    A & B & \\
    \hline
    True & True & Possible\\
    True & False & Not Possible\\
    False & True & Not Possible\\
    False & False & Possible\\

\end{tabular}
\end{center}

If And Only If is often abbreviated to ``iff''. It is also said that $A$ and $B$
are ``Equivalent''.
The symbol to denote this is $\iff$

\subsubsection*{And, Or and Not}
And, Or, and Not are also given specific meaning in proof writing:
These meanings can be shown using a truth table
\begin{center}
   \begin{tabular}{l|l||l}
       $A$ & $B$ & $A$ and $B$\\
       \hline
       True & True & True\\
       True & False & False\\
       False & True & False\\
       False & False & False\\
   \end{tabular}

\end{center}
 
\begin{center}
   \begin{tabular}{l||l}
       $A$ & not $A$\\
       \hline
       True & False\\
       False & True\\
   \end{tabular}
\end{center}

``And'' and ``Not'' are more similar to their English counterparts in use than
``Or''. Where the mathematical ``Or'' differentiates is in its acceptance of both.
If I say: ``I want ice cream or cookies'' it is not expected to see me wolfing both
down like a freaking hog.

\begin{center}
    \begin{tabular}{l|l||l}
        $A$ & $B$ & $A$ or $B$\\
        \hline
        True & True & True\\
        True & False & True\\
        False & True & True\\
        False & False & False\\
    \end{tabular}
\end{center}

Its like the || operator in C! Makes sense as I assume mathematical structure of
language influenced programming. 
    
\subsubsection*{Levels to This Splish}
Not all true mathematical statements deserve the prestige of being called a ``Theorem''.
$3+9=12$ isn't that special. Here are the alternative words used:
\begin{itemize}
    \item\textbf{Result:} Exudes modesty, can be significant or not. Both accepted
    \item\textbf{Fact:} Minor, $1+1=2$
    \item\textbf{Proposition:} A little more important than a fact
    \item\textbf{Lemma:} A smaller theorem used to build a larger proof.
    \item\textbf{Corollary:} A teeny little proof that is just an expansion on
        a previously proven theorem.
    \item\textbf{Claim:} Like a lemma. Used to build/organize a larger proof. Can be
        specific to the context of the enclosing proof.
\end{itemize}

\subsubsection*{Vacuous Truth}
A statement in which $A$ or $B$ is just freaking impossible.
\begin{stmt}
    \textbf{Vacuous:} If an integer is both a perfect square and a prime, then its negative 
\end{stmt}
While this statement is not useful as there is no possible number for which to apply it is 
still considered true. This is because no example disproves it. Vacuous is true but useless.

\newpage
\section*{Chapter 2}
\addcontentsline{toc}{section}{Chapter 2}
\end{document}

